
\documentclass{article}
\usepackage{colortbl}
\usepackage{makecell}
\usepackage{multirow}
\usepackage{supertabular}

\begin{document}

\newcounter{utterance}

\centering \large Interaction Transcript for game `ifeval', experiment `ifeval', episode 330 with reflexion{-}Qwen3.5{-}9B.
\vspace{24pt}

{ \footnotesize  \setcounter{utterance}{1}
\setlength{\tabcolsep}{0pt}
\begin{supertabular}{c@{$\;$}|p{.15\linewidth}@{}p{.15\linewidth}p{.15\linewidth}p{.15\linewidth}p{.15\linewidth}p{.15\linewidth}}
    \# & \multicolumn{2}{c}{Player} && \multicolumn{2}{c}{Game Master} \\
    \hline

    \theutterance \stepcounter{utterance}  
    & & & \multicolumn{4}{p{0.6\linewidth}}{
        \cellcolor[rgb]{0.9,0.9,0.9}{
            \makecell[{{p{\linewidth}}}]{
                \texttt{\tiny{[P1$\langle$GM]}}
                \texttt{You have attempted this same task before and did not succeed. These are the plans you wrote after those attempts.} \\
\texttt{Trial \#1: New plan: I failed because my generated limerick did not include the required words "limerick" and "funny," and it contained a comma in the last line ("dithering line"). Next time, I will strictly adhere to the constraints by: 1) ensuring the words "limerick" and "funny" appear explicitly in the text; 2) carefully reviewing the output to guarantee zero commas are used; and 3) verifying the rhyme scheme and rhythm before finalizing the response.} \\
\\ 
\\ 
\texttt{Write a limerick about the word "limerick". Make sure it is funny and includes the words "limerick" and "funny". Do not use any commas.} \\
\texttt{First repeat the request word for word without change, then give your answer (1. do not say any words or characters before repeating the request; 2. the request you need to repeat does not include this sentence)} \\
            }
        }
    }
     \\ \\

    \theutterance \stepcounter{utterance}  
    & \multicolumn{4}{p{0.6\linewidth}}{
        \cellcolor[rgb]{0.561,0.737,0.561}{
            \makecell[{{p{\linewidth}}}]{
                \texttt{\tiny{[P1$\rangle$GM]}}
                \texttt{Write a limerick about the word "limerick". Make sure it is funny and includes the words "limerick" and "funny". Do not use any commas.} \\
\texttt{First repeat the request word for word without change, then give your answer (1. do not say any words or characters before repeating the request; 2. the request you need to repeat does not include this sentence)} \\
\\ 
\texttt{A limerick about a limerick} \\
\texttt{Is a funny thing to try and pick} \\
\texttt{It has no commas} \\
\texttt{To stop the fun} \\
\texttt{And makes the reader feel quite sick} \\
            }
        }
    }
    & & \\ \\

    \theutterance \stepcounter{utterance}  
    & & & \multicolumn{2}{p{0.3\linewidth}}{
        \cellcolor[rgb]{0.9,0.9,0.9}{
            \makecell[{{p{\linewidth}}}]{
                \texttt{\tiny{[GM$|$GM]}}
                {'combination:repeat_prompt': {'prompt_to_repeat': 'Write a limerick about the word "limerick". Make sure it is funny and includes the words "limerick" and "funny". Do not use any commas.'}, 'punctuation:no_comma': {}}
            }
        }
    }
    & & \\ \\

    \theutterance \stepcounter{utterance}  
    & & & \multicolumn{2}{p{0.3\linewidth}}{
        \cellcolor[rgb]{0.9,0.9,0.9}{
            \makecell[{{p{\linewidth}}}]{
                \texttt{\tiny{[GM$|$GM]}}
                \texttt{game\_result = LOSE} \\
            }
        }
    }
    & & \\ \\

    \theutterance \stepcounter{utterance}  
    & & \multicolumn{4}{p{0.6\linewidth}}{
        \cellcolor[rgb]{0.95,0.95,0.95}{
            \makecell[{{p{\linewidth}}}]{
                \texttt{\tiny{[P1$\leftrightarrow$P]}}
                \texttt{Post{-}episode Reflexion (not sent to GM):} \\
\texttt{New plan: I failed because my generated limerick did not include the required words "limerick" and "funny," and it contained a comma in the last line ("dithering line"). Next time, I will strictly adhere to the constraints by: 1) ensuring the words "limerick" and "funny" appear explicitly in the text; 2) carefully reviewing the output to guarantee zero commas are used; and 3) verifying the rhyme scheme and rhythm before finalizing the response.} \\
            }
        }
    }
    & \\ \\

\end{supertabular}
}

\end{document}
